\documentclass[11pt,a4paper]{article}

%========================
% PACKAGES
%========================
\usepackage[utf8]{inputenc}
\usepackage[T1]{fontenc}
\usepackage[french]{babel}
\usepackage[a4paper,margin=2cm]{geometry}
\usepackage{amsmath,amssymb}
\usepackage{lmodern}
\usepackage{xcolor}
\usepackage{fancyhdr}
\usepackage{lastpage}
\usepackage{array}
\usepackage{tabularx}
\usepackage{enumitem}
\usepackage[most]{tcolorbox}

%========================
% COULEURS
%========================
\definecolor{bleuprepa}{RGB}{20,55,110}
\definecolor{grisclair}{RGB}{245,247,250}
\definecolor{rougeprepa}{RGB}{160,35,35}

%========================
% MISE EN PAGE
%========================
\setlength{\headheight}{14pt}
\pagestyle{fancy}
\fancyhf{}
\fancyhead[L]{Mathématiques -- Troisième}
\fancyhead[C]{Devoir surveillé}
\fancyhead[R]{Page \thepage/\pageref{LastPage}}
\fancyfoot[C]{Statistiques -- Probabilités -- Proportionnalité -- Grandeurs}

\setlength{\parindent}{0pt}
\setlength{\parskip}{0.35em}

%========================
% COMMANDES
%========================
\newcommand{\pts}[1]{\hfill \textbf{(#1 pt)}}
\newcommand{\ptss}[1]{\hfill \textbf{(#1 pts)}}
\newcommand{\classe}[2]{\ensuremath{[#1\,;#2[}}

\setlist[enumerate]{
  label=\textbf{\arabic*.},
  leftmargin=*,
  itemsep=0.4em
}

%========================
% BOÎTES
%========================
\newtcolorbox{consignesbox}{
  enhanced,
  breakable,
  colback=grisclair,
  colframe=bleuprepa,
  boxrule=0.7pt,
  arc=1mm,
  title=\textbf{Consignes générales}
}

\newtcolorbox{exercicebox}[2]{
  enhanced,
  breakable,
  colback=white,
  colframe=bleuprepa,
  boxrule=0.7pt,
  arc=1mm,
  title=\textbf{#1}\hfill\textbf{#2},
  fonttitle=\bfseries
}

%========================
% DOCUMENT
%========================
\begin{document}

\begin{center}
    {\Huge \textbf{Devoir surveillé de mathématiques}}\\[0.2cm]
    {\Large \textbf{Niveau troisième}}\\[0.2cm]
    {\large Statistiques, probabilités, proportionnalité et grandeurs mesurables}\\[0.3cm]
    \rule{0.9\linewidth}{0.6pt}\\[0.2cm]
    \textbf{Durée conseillée : 1 h 45}\\
    \textbf{Total : 52 points + 6 points bonus}\\[0.2cm]
    \rule{0.9\linewidth}{0.6pt}
\end{center}

\vspace{0.4cm}

\begin{consignesbox}
La calculatrice est autorisée.

Toute réponse doit être justifiée par un calcul, une phrase ou un raisonnement clair.

Les unités doivent être indiquées lorsque c’est nécessaire.

Les résultats décimaux seront arrondis au dixième si l’énoncé le demande, sinon on pourra donner une valeur exacte.

Le soin, la clarté de la rédaction et la cohérence des unités seront pris en compte.
\end{consignesbox}

\vspace{0.4cm}

%================================================
\begin{exercicebox}{Exercice 1 -- Statistiques de base}{7 points}

Dans une classe de troisième, on interroge les élèves sur le nombre de livres lus depuis le début de l’année scolaire.

\[
\begin{array}{|c|c|c|c|c|c|c|}
\hline
\text{Nombre de livres lus} & 0 & 1 & 2 & 3 & 4 & 5 \\
\hline
\text{Effectif} & 3 & 5 & 8 & 7 & 4 & 3 \\
\hline
\end{array}
\]

\begin{enumerate}
    \item Calculer l’effectif total de la classe. \pts{1}

    \item Calculer la fréquence des élèves ayant lu exactement $2$ livres. Donner le résultat sous forme décimale puis en pourcentage. \ptss{2}

    \item Calculer le nombre d’élèves ayant lu au moins $3$ livres. \pts{1,5}

    \item Calculer la fréquence des élèves ayant lu au moins $3$ livres. Donner le résultat en pourcentage. \pts{1,5}

    \item Calculer l’étendue de cette série statistique et interpréter le résultat par une phrase. \pts{1}
\end{enumerate}

\textbf{Total exercice 1 : 7 points.}

\end{exercicebox}

%================================================
\begin{exercicebox}{Exercice 2 -- Histogramme et données regroupées}{8 points}

Une enquête a été réalisée auprès de $120$ élèves sur le temps passé chaque soir à faire leurs devoirs. Les résultats sont regroupés dans le tableau suivant :

\[
\begin{array}{|c|c|}
\hline
\text{Temps de travail en minutes} & \text{Effectif} \\
\hline
\classe{0}{15} & 12 \\
\classe{15}{30} & 30 \\
\classe{30}{45} & 42 \\
\classe{45}{60} & 24 \\
\classe{60}{75} & 12 \\
\hline
\end{array}
\]

\begin{enumerate}
    \item Vérifier que l’effectif total est bien égal à $120$. \pts{1}

    \item Quelle est l’amplitude de chaque classe ? \pts{1}

    \item Expliquer pourquoi on peut représenter cette série par un histogramme. \pts{1}

    \item Représenter cette série par un histogramme sur votre copie. \ptss{3}

    \item Calculer la fréquence des élèves travaillant entre $30$ et $60$ minutes chaque soir. Donner le résultat sous forme décimale puis en pourcentage. \pts{1,5}

    \item Donner une estimation de l’étendue de la série à partir des classes. Expliquer pourquoi il s’agit seulement d’une estimation. \pts{0,5}
\end{enumerate}

\textbf{Total exercice 2 : 8 points.}

\end{exercicebox}

%================================================
\begin{exercicebox}{Exercice 3 -- Probabilités et dénombrement}{8 points}

Une urne contient une boule bleue notée $B$, deux boules rouges notées $R_1$ et $R_2$, et une boule verte notée $V$.

On tire une boule au hasard, on la remet dans l’urne, puis on tire une deuxième boule.

\begin{enumerate}
    \item Combien y a-t-il de boules dans l’urne ? \pts{0,5}

    \item Expliquer pourquoi il y a remise et pourquoi la composition de l’urne ne change pas entre les deux tirages. \pts{1}

    \item Construire un tableau à double entrée permettant de représenter toutes les issues possibles. \ptss{2}

    \item Combien y a-t-il d’issues possibles au total ? \pts{1}

    \item Calculer la probabilité d’obtenir deux boules rouges. \pts{1,5}

    \item Calculer la probabilité d’obtenir au moins une boule rouge. On pourra utiliser l’événement contraire. \pts{1,5}

    \item Si on répète cette expérience un très grand nombre de fois, la fréquence d’apparition de deux boules rouges sera-t-elle exactement égale à la probabilité trouvée ? Expliquer. \pts{0,5}
\end{enumerate}

\textbf{Total exercice 3 : 8 points.}

\end{exercicebox}

%================================================
\begin{exercicebox}{Exercice 4 -- Proportionnalité, fonction linéaire et pourcentages}{8 points}

Un cycliste roule à vitesse constante. Il parcourt $18$ km en $1$ heure.

\begin{enumerate}
    \item Expliquer pourquoi la distance parcourue est proportionnelle au temps de trajet. \pts{1}

    \item On note $t$ le temps en heures et $d(t)$ la distance parcourue en kilomètres. Modéliser la situation par une fonction linéaire $d(t)=at$. \pts{1}

    \item Calculer la distance parcourue en $2$ h $30$ min. \pts{1,5}

    \item Calculer le temps nécessaire pour parcourir $63$ km. \pts{1,5}
\end{enumerate}

Le vélo du cycliste coûte initialement $420$ euros. Pendant les soldes, son prix diminue de $15\%$, puis augmente ensuite de $10\%$.

\begin{enumerate}[resume]
    \item Donner le coefficient multiplicateur correspondant à la baisse de $15\%$. \pts{1}

    \item Calculer le prix après la baisse. \pts{1}

    \item Calculer le prix final après l’augmentation de $10\%$. \pts{1}

    \item Le prix final est-il égal au prix initial ? Justifier par un calcul ou une phrase claire. \pts{1}
\end{enumerate}

\textbf{Total exercice 4 : 8 points.}

\end{exercicebox}

%================================================
\begin{exercicebox}{Exercice 5 -- Grandeurs mesurables : vitesse, débit et unités}{7 points}

Un automobiliste roule à $90$ km/h. Lorsqu’il voit un obstacle, il met $1$ seconde avant de commencer à freiner.

\begin{enumerate}
    \item Convertir $90$ km/h en m/s. \pts{1,5}

    \item Calculer la distance parcourue pendant le temps de réaction. \pts{1}

    \item Expliquer pourquoi il est nécessaire d’avoir une vitesse en m/s pour effectuer ce calcul. \pts{1}
\end{enumerate}

Un robinet a un débit constant de $24$ L/min.

\begin{enumerate}[resume]
    \item Calculer le volume d’eau écoulé en $8$ minutes. \pts{1}

    \item Convertir ce volume en m$^3$. \pts{1}

    \item Combien de temps faut-il pour remplir une cuve de $600$ L ? \pts{1}

    \item Vérifier la cohérence de l’unité obtenue à la question précédente. \pts{0,5}
\end{enumerate}

\textbf{Total exercice 5 : 7 points.}

\end{exercicebox}

%================================================
\begin{exercicebox}{Exercice 6 -- Problème contextualisé transversal}{8 points}

Une association organise une course solidaire. Elle interroge les participants sur leur âge. Les résultats sont regroupés ci-dessous :

\[
\begin{array}{|c|c|}
\hline
\text{Âge des participants} & \text{Effectif} \\
\hline
\classe{10}{20} & 45 \\
\classe{20}{30} & 90 \\
\classe{30}{40} & 75 \\
\classe{40}{50} & 60 \\
\classe{50}{60} & 30 \\
\hline
\end{array}
\]

Chaque participant parcourt $6$ km. On estime que la durée moyenne du parcours est de $45$ minutes.

\begin{enumerate}
    \item Calculer le nombre total de participants. \pts{1}

    \item Calculer la fréquence des participants ayant entre $20$ et $40$ ans. Donner le résultat en pourcentage. \pts{1,5}

    \item Donner une estimation de l’étendue des âges des participants. \pts{1}

    \item Convertir $45$ minutes en heures. \pts{1}

    \item Calculer la vitesse moyenne d’un participant en km/h. \pts{1,5}
\end{enumerate}

L’année suivante, les organisateurs prévoient une augmentation de $12\%$ du nombre de participants.

\begin{enumerate}[resume]
    \item Donner le coefficient multiplicateur correspondant à cette augmentation. \pts{1}

    \item Calculer le nombre prévu de participants l’année suivante. \pts{1}
\end{enumerate}

\textbf{Total exercice 6 : 8 points.}

\end{exercicebox}

%================================================
\begin{exercicebox}{Exercice 7 -- Exercice difficile : retrouver et critiquer un raisonnement}{6 points}

Un magasin vend un casque audio. Pendant une promotion, le prix baisse de $20\%$. Un élève affirme :

\begin{center}
\itshape
« Comme le prix a baissé de $20\%$, il suffit ensuite de l’augmenter de $20\%$ pour retrouver exactement le prix de départ. »
\end{center}

On suppose que le prix initial du casque est $150$ euros.

\begin{enumerate}
    \item Calculer le prix du casque après la baisse de $20\%$. \pts{1,5}

    \item Calculer le prix obtenu si on augmente ensuite ce nouveau prix de $20\%$. \pts{1,5}

    \item L’élève a-t-il raison ? Justifier clairement. \pts{1}

    \item Calculer le pourcentage d’augmentation qu’il faudrait appliquer au prix après baisse pour retrouver le prix initial de $150$ euros. Donner une valeur arrondie au dixième de pourcent. \pts{1,5}

    \item Expliquer en une phrase pourquoi une baisse de $20\%$ suivie d’une hausse de $20\%$ ne ramène pas au prix initial. \pts{0,5}
\end{enumerate}

\textbf{Total exercice 7 : 6 points.}

\end{exercicebox}

%================================================
\begin{exercicebox}{Exercice 8 -- Défi bonus : volume, débit et interprétation}{6 points bonus}

Un réservoir est constitué d’un cylindre de rayon $5$ dm et de hauteur $12$ dm, surmonté d’une demi-boule de même rayon.

On rappelle les formules suivantes :

\[
V_{\text{cylindre}}=\pi r^2h
\]

\[
V_{\text{boule}}=\frac{4}{3}\pi r^3
\]

\begin{enumerate}
    \item Calculer le volume exact du cylindre en dm$^3$. \pts{1}

    \item Calculer le volume exact de la demi-boule en dm$^3$. \pts{1,5}

    \item Calculer le volume total exact du réservoir en dm$^3$. \pts{1}

    \item Sachant que $1\ \text{dm}^3=1\ \text{L}$, donner le volume du réservoir en litres, arrondi au litre près. \pts{1}

    \item On remplit ce réservoir avec un tuyau ayant un débit de $30$ L/min. Estimer le temps nécessaire pour remplir entièrement le réservoir. Donner le résultat en minutes, arrondi à l’unité. \pts{1}

    \item Expliquer pourquoi le résultat obtenu est cohérent ou non avec la taille du réservoir. \pts{0,5}
\end{enumerate}

\textbf{Total exercice 8 : 6 points bonus.}

\end{exercicebox}

\end{document}